\section{Introduction}
\label{sec:introduction}

Written language varies across regions, communities, and media. Search,
translation, and accessibility systems nevertheless often expect a single
standardized form. Dialect normalization---rewriting a non-standard input into
that form without changing its content---can bridge this mismatch, but it is
especially difficult for the communities that have the least labeled data.
The same spelling may be a dialect marker in one context and a person or place
name in another; aggressive correction therefore trades recall for semantic
damage.

Multilingual pretraining transfers useful abstractions across languages
\citep{devlin2019bert,conneau2020xlmr}, and byte-level models reduce reliance
on brittle tokenization \citep{xue2022byt5}. Yet neither mechanism gives a
model explicit access to the few verified transformations that a community
has already annotated. Retrieval augmentation can supply such evidence at
inference time \citep{lewis2020rag,khandelwal2021knnmt}, but naive retrieval
creates a new failure mode: a superficially similar example can override a
correct model prediction.

We address this tension with \method, shown in Figure~\ref{fig:architecture}.
The system treats retrieval as a selective intervention. A byte-level encoder--
decoder first produces a distribution and an uncertainty estimate. Only
uncertain spans query a compact memory of verified normalization pairs. A gate
then combines generation and memory scores using retrieval relevance, model
uncertainty, and named-entity compatibility. Every copied form remains linked
to the example that supported it, providing a lightweight audit trail.

This illustrative study makes three contributions:
\begin{itemize}
  \item an uncertainty-triggered retrieval policy that avoids needless memory
        lookups for confident spans;
  \item an evidence gate that jointly scores relevance, local agreement, and
        entity preservation; and
  \item a controlled evaluation protocol covering accuracy, calibration,
        robustness, and harm-relevant entity corruption.
\end{itemize}

% -----------------------------------------------------------------------------
